% Evaluation writeup: planner comparison across environments and horizons.
% Numbers are filled in from the sweep by writeup/fill_tables.py, which reads
% the seed_success_log.txt files directly -- do not hand-edit the RESULTS
% tables, edit the generator.
\documentclass[11pt]{article}
\usepackage[margin=1in]{geometry}
\usepackage{booktabs}
\usepackage{amsmath}
\usepackage{graphicx}
\usepackage{xcolor}
\usepackage{multirow}
\usepackage[colorlinks=true,linkcolor=blue,citecolor=blue]{hyperref}

\newcommand{\TODO}[1]{\textcolor{red}{[#1]}}
\newcommand{\na}{--}

\title{Gradient-based planning with learned world models:\\
a controlled comparison across environments and horizons}
\author{}
\date{\today}

\begin{document}
\maketitle

\section{Protocol}

\paragraph{Seeding (important).} Prior to 2026-08-31 every run in this
repository used a degenerate seeding path: \texttt{plan.py} computed
\[
  \texttt{eval\_seed} = [\,\texttt{seed} \cdot n + 1 \;\text{for}\; n \in \operatorname{range}(\texttt{n\_evals})\,],
\]
which for \texttt{n\_evals=1} evaluates to $[1]$ for \emph{every} seed, so all
runs planned the same start/goal pair. In addition \texttt{utils.seed()} was
called only when \texttt{planner\_seed} was non-null, which it never was, so
the random action initialisation was drawn from an unseeded generator. Both
were fixed; all numbers in this document come from the corrected code, and
earlier numbers in the repository are not comparable to them.

\paragraph{Success criterion.} Success is environment-defined. For point maze
(\texttt{env/pointmaze/point\_maze\_wrapper.py}) it is
$\lVert g_{1:2} - s_{1:2}\rVert < 0.5$, i.e.\ a threshold on \emph{position
only}; the quantity logged as \texttt{state\_dist} is the full state norm
including velocity and does \emph{not} determine success. Where we report a
distance we report the success-determining one, minimised over the
optimisation trace (so it does not confound ``how close did it get'' with
``when did it stop'').

\paragraph{Goal source.} All runs use \texttt{goal\_source=dset}. Episodes are
$100$ frames (point maze) and $201$ frames (reacher) with frameskip $5$, so
\texttt{goal\_source=dset} admits $\texttt{goal\_H} \le 19$ and $\le 40$
respectively; the horizons studied here all fit. Under \texttt{dset} the goal
lies \texttt{goal\_H} steps ahead on a real trajectory, so horizon scales task
difficulty; under \texttt{random\_state} the goal is arbitrary and
\texttt{goal\_H} changes only the planner's action horizon.

\paragraph{Budget.} All planners are given the same optimisation-step budget
of $150$ steps. Step budget and \emph{compute} budget are very different here
(Table~\ref{tab:compute}), so each table states which it holds fixed.

\section{Environments and models}

\begin{table}[h]
\centering
\begin{tabular}{llll}
\toprule
Environment & World model & Latent & Status \\
\midrule
Reacher    & LeWM (JEPA)  & visual, $192$-d single token & running \\
Cube       & LeWM (JEPA)  & visual                        & \textbf{no data} \\
Point maze & DINO-WM      & proprio, $10$-d               & running \\
Push-T     & DINO-WM      & visual, $196$ ViT patches     & queued \\
Wall       & DINO-WM      & visual                        & \textbf{no data} \\
\bottomrule
\end{tabular}
\caption{Environments. Point maze and Push-T carry $196$ ViT patch tokens per
frame, costing ${\sim}2.1$\,GB per retained rollout, which bounds population
sizes on a $16$\,GB card.}
\end{table}

\section{Planners and compute accounting}

\begin{table}[h]
\centering
\begin{tabular}{lllr}
\toprule
Planner & Mechanism & Configuration & Rollouts/step \\
\midrule
GD    & Adam on actions                & \texttt{num\_samples}$=1$                      & $1$ \\
GRASP & lifted virtual states + penalty & \texttt{n\_parallel}$=1$                       & $1$ \\
CEM   & sampling + elite refit          & $30$ samples $\times$ $5$ iters, top-$3$        & $30$ \\
\bottomrule
\end{tabular}
\caption{Per-step world-model rollouts. Stock \texttt{cem.yaml} uses
\texttt{num\_samples}$=300$, \texttt{opt\_steps}$=30$ ($9{,}000$ rollouts)
against GD's $150$, so it must be retuned to be comparable; see
Section~\ref{sec:cem}.}
\label{tab:compute}
\end{table}

\subsection{Tuning CEM to the budget}
\label{sec:cem}
CEM is matched to GD on \emph{world-model rollouts}, not on optimiser
iterations: \texttt{num\_samples}$=30$, \texttt{opt\_steps}$=5$,
\texttt{topk}$=3$, giving $30 \times 5 = 150$ rollouts against GD's
$150 \times 1$. Stock \texttt{cem.yaml} ($300 \times 30 = 9{,}000$) would
otherwise give CEM $60\times$ the compute.

Rollout-matching and wall-clock-matching are not the same experiment here,
because CEM batches its population into one GPU call while GD is
launch-overhead bound; an earlier reacher measurement found near-identical
wall-clock at a $247\times$ difference in model calls. We therefore report
per-run wall clock alongside every success rate. On the reacher probe the two
happen to agree closely (CEM $101$\,s vs GD $123$\,s at equal rollouts), but
this need not hold on Push-T, where $196$ patch tokens per frame make batched
evaluation relatively cheaper, so the wall-clock column should be read
before drawing a compute-fairness conclusion per environment.

\section{Results}

\begin{table}[h]
\centering
\begin{tabular}{llccc}
\toprule
Environment & $H$ & GD & GRASP & CEM \\
\midrule
\multirow{3}{*}{Reacher}
 & 5  & \TODO{} & \TODO{} & \TODO{} \\
 & 10 & \TODO{} & \TODO{} & \TODO{} \\
 & 15 & \TODO{} & \TODO{} & \TODO{} \\
\midrule
\multirow{3}{*}{Cube}
 & 5  & \TODO{} & \TODO{} & \TODO{} \\
 & 10 & \TODO{} & \TODO{} & \TODO{} \\
 & 15 & \TODO{} & \TODO{} & \TODO{} \\
\midrule
\multirow{3}{*}{Point maze}
 & 5  & \TODO{} & \TODO{} & \TODO{} \\
 & 10 & \TODO{} & \TODO{} & \TODO{} \\
 & 15 & \TODO{} & \TODO{} & \TODO{} \\
\midrule
\multirow{3}{*}{Push-T}
 & 5  & \TODO{} & \TODO{} & \TODO{} \\
 & 10 & \TODO{} & \TODO{} & \TODO{} \\
 & 15 & \TODO{} & \TODO{} & \TODO{} \\
\midrule
\multirow{3}{*}{Wall}
 & 5  & \TODO{} & \TODO{} & \TODO{} \\
 & 10 & \TODO{} & \TODO{} & \TODO{} \\
 & 15 & \TODO{} & \TODO{} & \TODO{} \\
\bottomrule
\end{tabular}
\caption{Success rate (\%) over $100$ seeds, $150$ optimisation steps.
Seeds index distinct start/goal pairs (see Protocol).}
\end{table}

\section{Data availability}

Three of the five environments required work before they could be evaluated,
and two could not be recovered at all.

\begin{itemize}
  \item \textbf{Push-T} shipped as a $2.8$\,GB archive. Planning draws goals
        only from the validation split (\texttt{plan.py} selects
        \texttt{dset["valid"]}), so only the $21$ validation episodes needed
        frame extraction ($2{,}514$ PNGs, $22$\,MB) rather than all $18{,}685$
        training videos. Validation episodes run $76$--$167$ frames, so at
        frameskip $5$ the $H{=}15$ condition ($76$ frames) is satisfied by
        every episode, but only just.
  \item \textbf{Cube} is unavailable: the local
        \texttt{ogbench/lewm\_cube\_raw} tree contains no dataset, only an
        interrupted HuggingFace download (a single \texttt{.incomplete} blob).
        \texttt{conf/env/cube.yaml} expects
        \texttt{ogbench/cube\_single\_expert.h5}, which is absent.
  \item \textbf{Wall} has no local data and its host was unreachable.
\end{itemize}

\paragraph{Environment deviations.} \texttt{pymunk}, \texttt{pygame} and
\texttt{shapely} were absent and required for the Push-T simulator.
\texttt{shapely} is installed at $2.1.2$ rather than the pinned $2.0.3$,
because the pinned build is compiled against NumPy 1.x and this environment
runs NumPy $2.2.6$.

\paragraph{Code fixes required for the sweep.} \texttt{GDPlanner} loaded a
replay buffer and built a $k$-NN graph unconditionally, including at
\texttt{hjb\_tau}$=0$ where $V$ is multiplied by zero; this is now skipped
(verified to reproduce prior results exactly). \texttt{planning/cem.py} wrote
its log to a hardcoded cluster path and crashed after planning had already
succeeded. \texttt{planning/grasp.py} writes no
\texttt{seed\_success\_log.txt} at all, so results are harvested uniformly
from the evaluator's own output rather than from that file.

\section{Notes and caveats}
\begin{itemize}
  \item \textbf{Point maze data.} \texttt{states.pth}/\texttt{actions.pth}
        cover $2{,}000$ episodes but only $200$ rendered observation sets are
        present on disk, so the effective dataset is ${\sim}10\%$ of nominal.
        The derived latent buffer holds $3{,}800$ transitions against
        reacher's $55{,}000$.
  \item \textbf{Baseline strength.} Stock \texttt{gd.yaml} uses
        \texttt{sample\_type=randn} and \texttt{action\_reg=0}. On point maze
        at $H{=}10$, switching to \texttt{sample\_type=zero} with
        \texttt{action\_reg}$=0.1$ moved GD from $0/4$ to $3/4$ on a selected
        subset, so the stock configuration is a weak baseline and should be
        reported as one configuration among several rather than as ``GD''.
  \item \textbf{GRASP is not byte-identical to the published version} --- it
        carries a \texttt{\_best\_parallel\_indices()} patch and a
        \texttt{rollout\_encoded} addition for LeWM. Both make it stronger.
\end{itemize}

\end{document}
